\subsection{Natural Numbers and Whole Numbers} 
There are several standard sets of numbers we discuss in mathematics. We build them up formally in the same way we built them up informally as we first learned about numbers.
\begin{definition}[Natural Numbers]
The \term{natural numbers} (sometimes called the \term{counting numbers}) are the set
\[
\mathbb{N}=\{1,2,3,\dots\}
\]
\end{definition}
\begin{itemize}
\item The natural numbers are all \makeblank
\item They do not have any \makeblank parts
\end{itemize}
The natural numbers are the numbers that give an intuitive answer to ``how many?'' questions. We can add in one more (less intuitive) answer to such a question, and our math gets much more powerful.
\begin{definition}[Whole Numbers]
The \term{whole numbers} are the set
\[
\mathbb{W}=\{0,1,2,3,\dots\}
\]
The difference from the natural numbers is the number \makeblanksmall
\end{definition}
\begin{Note}
The number $0$ is \emph{not} considered a \makeblank number. (It isn't a \makeblank number, either---it's in a cateogry of its own.)
\end{Note}
\begin{Example}
Which of these numbers are natural numbers? Whole numbers?
\begin{itemize}
\item $-2$ \checkbox{Natural Number} \checkbox{Whole Number}
\makeblanknew
\item $0$ \checkbox{Natural Number} \checkbox{Whole Number}
\makeblanknew
\item $3$ \checkbox{Natural Number} \checkbox{Whole Number}
\makeblanknew
\item $\frac{10}{5}$ \checkbox{Natural Number} \checkbox{Whole Number}
\makeblanknew
\item $\frac{11}{5}$ \checkbox{Natural Number} \checkbox{Whole Number}
\makeblanknew
\end{itemize}
\end{Example}